% Chapter 1

\chapter{Specifica dei requisiti} % Main chapter title

\label{Chapter1} % For referencing the chapter elsewhere, use \ref{Chapter1} 

%----------------------------------------------------------------------------------------

% Define some commands to keep the formatting separated from the content 
% \newcommand{\keyword}[1]{\textbf{#1}}
% \newcommand{\tabhead}[1]{\textbf{#1}}
% \newcommand{\code}[1]{\texttt{#1}}
% \newcommand{\file}[1]{\texttt{\bfseries#1}}
% \newcommand{\option}[1]{\texttt{\itshape#1}}

%----------------------------------------------------------------------------------------

% =========================================================== %
%                 Section: General Information                %
% =========================================================== %
\section{Scopo del documento dei requisiti}
\label{chap1:sec1:general_information}
Il documento \textit{SRS}, riportato nei primi tre capitoli, si propone di accogliere in forma organica una descrizione di tutte le specifiche tecniche e funzionali che caratterizzano il sistema per la gestione di issue in progetti software. \\
Esso è stato redatto seguendo le indicazioni dello standard IEEE/ANSI 830-1998 (IEEE 1998) per i documenti dei requisiti.

% =========================================================== %
%      Section: General Requirements of Graduate School       %
% =========================================================== %
\section{Scopo del prodotto}

Il prodotto software che si vuole realizzare prende il nome di BugBoard26, una piattaforma che consentirà a un team di sviluppo di segnalare problemi relativi a un progetto, monitorare lo stato, assegnarli a membri del team e tenere traccia delle attività di risoluzione. Il sistema deve fornire le seguenti funzionalità:

\begin{enumerate}
  \item Il sistema deve fornire un meccanismo di autenticazione semplice e sicuro, basato su email e
        password. Le informazioni gestite dall’applicazione sono critiche per l’azienda, ed è fondamentale preservarne l’integrità e la segretezza. Il sistema viene fornito con un account da amministratore già attivo, con credenziali di default. Un amministratore può creare ulteriori utenze, specificando una email, una password, e indicando se quell’utenza sarà "normale" oppure "di amministrazione".
  \item Tutti gli utenti autenticati possono segnalare una issue indicando almeno un titolo e una
        descrizione. Alcuni utenti potrebbero voler specificare anche una priorità e sarebbe gradita la possibilità di allegare un’immagine. Le issue possono essere di diverso tipo: question (per richieste di chiarimenti), bug (per segnalare malfunzionamenti), documentation (per segnalare problemi relativi alla documentazione), e feature (per indicare la richiesta o il suggerimento di nuove funzionalità). Le issue create sono inizialmente nello stato "todo".
  \item Il sistema deve offrire una vista riepilogativa delle issue, con la possibilità di filtrare        o ordinare i risultati in base a criteri come tipologia, stato, priorità o altri parametri         rilevanti.
  \item Ogni bug dovrebbe poter essere assegnato (da un amministratore) a un membro del team. Quando
        viene assegnato un bug a un utente, quest’ultimo riceve una notifica. Gli utenti possono
        visualizzare i bug loro assegnati.
  \item È necessaria una funzione di ricerca che consenta di trovare bug in base a parole chiave.
  \item Il sistema dovrebbe suggerire automaticamente a chi assegnare un bug, basandosi sul carico di
        lavoro corrente degli utenti.
\end{enumerate}

\section{Riferimenti}
Qui di seguito sono elencati i documenti e gli indirizzi web a cui fa riferimento la presente specifica dei requisiti, nonché i 
link al codice sorgente e al repository Docker
\begin{itemize}
    \item Slide del corso di Ingegneria del software 2025/2026 fornite dai professori P.Tramontana e B.Breve.
    \item Ian, Sommerville. \textit{Ingegneria del software}. 8\textsuperscript{a} edizione, Pearson Paravia Bruno Mondadori S.p.A., 2007.
    \item Modello \href{https://ieeexplore.ieee.org/document/720574}{IEEE/ANSI 830-1998} per strutturare il documento dei requisiti.
    \item Enciclopedia online \href{https://it.wikipedia.org/wiki/Pagina_principale}{Wikipedia}. 
    \item Codice sorgente \href{https://github.com/ubiestubi/ProgettoIngSoft}{Back-end}.
    \item Codice sorgente \href{https://github.com/GiandomenicoIameo/BugBoard26/tree/main}{Front-end}.
    \item Link alla \href{https://communitystudentiunina-my.sharepoint.com/:f:/g/personal/u_musto_studenti_unina_it/IgArB_oG9uFMQ5LTcVQUakrmAcshTORCLn3K1V4J0z0thnY?e=ka4gLm}{directory} contenente l'immagine Docker e il file compose per il back-end
\end{itemize}

\section{Descrizione del resto del documento}
La restante parte del presente documento è suddivisa in due ulteriori capitoli che forniscono una descrizione sempre più dettagliata delle funzionalità del prodotto software in questione.
Il capitolo immediatamente successivo, \textit{Descrizione generale}, è composto dai seguenti paragrafi:

\begin{enumerate}
  \item \textit{Prospettiva del prodotto:} descrive il contesto o ambiente operativo in cui sarà utilizzato il prodotto software in questione. Verrà specificata l'architettura generale, se il prodotto è un membro di una famiglia di prodotti, un sostituto di altri sistemi oppure un prodotto completamente autonomo.
  \item \textit{Funzioni del prodotto:} illustra i principali servizi (o funzionalità) del prodotto 
  corredati da diagrammi. Il livello di dettaglio usato in questo paragrafo è a un alto livello; i dettagli delle funzionalità saranno forniti del capitolo 3.
  \item \textit{Caratteristiche utente:} esplicità le conoscenze richieste per poter utilizzare il prodotto in questione e fornisce una descrizione generale degli utilizzatori previsti dal sistema.
  \item \textit{Presupposti e dipendenze:} illustra le condizioni che il sistema presume siano verificate al momento del suo avvio, durante la sua esecuziione e dopo la terminazione. Verranno specificati anche \textit{elementi esterni} (altri sistemi, persone o componenti) da cui il prodotto dipende e che sono necessari per il suo corretto funzionamento.
\end{enumerate}

Nel terzo capitolo, \textit{Requisiti specifici}, verranno illustrate le funzionalità che il prodotto software realizza. Il capitolo è composto dai seguenti paragrafi:

\begin{enumerate}
  \item \textit{Requisiti funzionali:} fornisce una descrizione dettagliata dei casi d'uso che verranno presentati nel capitolo 2. Tali requisiti indicano come il sistema dovrebbe reagire a particolari input o come dovrebbe comportarsi in particolari situazioni.
  \item \textit{Requisiti non funzionali:} esplicita i vincoli sui servizi o sulle funzioni offerti dal sistema. In questa sezione, i requisiti imposti sono stati divisi secondo la \textit{classificazione di Sommerville}.
\end{enumerate}







